\section*{結言}
この日の観測では、モデルそのものの発表に加え、公開後の評価と実際に使える場所の拡大が大きな比重を占めた。Muse Spark OpenやNemotronの話題は前日の発表を引き継ぎ、Beamは配布・提供の経路を、Glimpseは視覚的な推論の方向性を示した。モデルの性能数値と実運用での価値をつなぐには、独立評価と再現可能な利用条件の整理が欠かせない。

Claudeを用いた暗号曲線の由来の解析は、生成AIが研究上の仮説検討に使われる事例として注目された。Claude CodeやCodexの改善は、日常の開発作業で待ち時間と手作業を減らす方向を示す。一方、中国の規制案に関する投稿のように一次資料が未確認の話題も含まれる。観測窓外の背景発表と当日の反応を混同せず、速報としての位置付けを明確にする必要がある。
